\subsection{Einzelbaumsegmentierung}

Ausgehend von den vorverarbeiteten Datensätzen wird im Folgenden die Einzelbaumsegmentierung durchgeführt. Dafür wird ein rasterbasierter Segmentierungsansatz gewählt. Da urbane Bäume oft eine hohe Heterogenität aufweisen, besteht bei punktwolkenbasierten Verfahren die Gefahr, dass benachbarte, ineinandergreifende Baumkronen fälschlicherweise miteinander verschmolzen werden. Da Straßenbäume in der Regel isoliert stehen und selten mehrschichtige Waldstrukturen bilden, ist der rasterbasierte Informationsverlust vertretbar. Einschränkungen der Segmentierungsgenauigkeit treten hierbei meist nur in dicht bewachsenen Parkanlagen auf. Bei der Skalierung auf ein gesamtes Stadtgebiet erweist sich ein rasterbasierter Ansatz als deutlich effizienter, zumal er auch bei moderater Punktdichte stabile Ergebnisse liefert \cite{munzingerMappingUrbanForest2022}. Die Umsetzung gliedert sich in die CHM-Erstellung, die Baumspitzendetektion sowie die final Segmentierung.